\documentclass{article}
\usepackage[utf8]{inputenc}
\usepackage{geometry}
\usepackage{enumitem}
\usepackage{hyperref}
\usepackage{xcolor}
\usepackage{graphicx}
\usepackage{array}
\usepackage{multirow}
\usepackage{amsmath, amssymb}
\usepackage{chemfig}
\usepackage{mhchem}
\usepackage{tikz}
\usepackage{setspace}
\usepackage{soul}
\usepackage{mathtools}
\usepackage{booktabs}
\usepackage{chemformula}
\usepackage{siunitx}
\usetikzlibrary{positioning, calc}
\usepackage{longtable}
\geometry{margin=1in}
\setlength{\parskip}{0.5em}
\usepackage{cancel}

\begin{document}

\begin{center}
    {\LARGE\textbf{SI Session Planning Form}}
\end{center}

\
\noindent
\begin{flushleft}
\textbf{SI Leader:} Brandon Cobb\\
\textbf{Session Week:} 13\\
\textbf{Session\#:} 13\\
\textbf{Instructor:} Professor Tramontozzi\\
\textbf{Old Plans:}\href{https://macombedu.sharepoint.com/:b:/s/ASC-SupplementalInstructionEmbeddedTutoring/EeSYqCwT_KhMolgJvXbVEjgBOyF_DevvS8BgCm1zPkQQJg?e=QJTt4o}{SharePoint}\\
\textbf{Session Goal:} Students will review and apply concepts of concentration and solution chemistry, including calculating molarity, mass percent, volume percent, and mass/volume percent. They will apply dilution principles and use solubility rules to determine whether a solution is unsaturated, saturated, or supersaturated, and to predict precipitation. Students will analyze factors affecting solubility, including temperature and the polarity of solute and solvent. They will also explore colligative properties such as freezing point depression and boiling point elevation for electrolytes and nonelectrolytes. Additionally, students will practice stoichiometric calculations, identify limiting reactants and theoretical yields, and classify chemical reactions including synthesis, decomposition, single and double replacement, combustion, neutralization, and redox reactions.

\end{flushleft}

\vspace{1em}
\section*{\underline{Opener}}
{\centering\large\textbf{\hl{Take Attendance}}\par}\vspace{-0.25\baselineskip}
\noindent
\begin{tabular}{|p{4.2cm}|p{9.8cm}|}
\hline
\textbf{Collaborative Learning Techniques} & Group Discussion\\
\hline
\textbf{Learning Strategy} & Lecture Review\\
\hline
\textbf{Time} & 12:00---12:10\\
\hline
\textbf{Materials} & Notes, homework, whiteboard, markers, ready to speak, periodic table \\
\hline
\textbf{Lecture/Textbook/Old Plans/Slide References} & Class roster, Chapter 
8: Solutions and Chapter 9: Chemical Reactions\\
\hline
\end{tabular}


\vspace{0.5cm}
{\large\textbf{Definitions}}
\\
\begin{tabular}{|p{0.9\linewidth}|}
\hline
\begin{minipage}[t]{0.9\linewidth}

\textbf{Group Discussion:} One person in the group is asked to present on a topic or review material for the group and then lead the discussion. This person should not be the regular group leader. The goal is to actively engage all members in sharing ideas, asking questions, and clarifying understanding through peer-to-peer interaction.

\textbf{KWLH Framework:}
\begin{itemize}
    \item \textbf{K (Know):} Students identify what they already know about the topic before the discussion begins.
    \item \textbf{W (Want to know):} Students list questions or areas of curiosity they want to explore during the discussion.
    \item \textbf{L (Learned):} During or after the discussion, students summarize what they have learned, including clarifications, new concepts, or insights from peers.
    \item \textbf{H (How):} Students reflect on how they can apply the knowledge gained or strategies used during the discussion to future learning or problem-solving.
\end{itemize}

\end{minipage} \\
\hline
\end{tabular}
\newpage
\noindent
{\large\textbf{Instructions (please provide a\hl{DETAILED} activity breakdown below (and/or attached}}\\

\begin{tabular}{|p{0.9\linewidth}|}
\hline
\vspace{0.2cm}
\textbf{KWLH:}
\begin{enumerate}
    \item How did summarizing what you already knew at the start of the session help you understand the lecture material?  
    \item Did predicting questions or directions for future lectures help focus your learning? How?  
    \item How did using your notes and handouts to find support for key concepts improve your understanding?  
    \item Did summarizing the main ideas at the end of the session help reinforce what you learned? In what way?  
    \item How helpful was writing a short paragraph and listing new vocabulary for consolidating your learning?  
    \item Did formulating potential exam questions from the session help you think critically about the material?  
    \item How did the KWLH activity help you connect previous knowledge with new concepts from the lecture?  
    \item What part of the KWLH process was most useful for your learning, and why?  
\end{enumerate} \\
\hline
\end{tabular}
\newpage

\noindent
{\large\textbf{Checking for Understanding (questions \& answers)}}\\

\begin{tabular}{|p{0.9\linewidth}|}
\hline
\begin{itemize}
\item How does adding a nonvolatile solute affect the freezing point of a solvent?
\item Why does dissolving an ionic compound usually lower the freezing point more than a covalent one?
\item What factors determine the magnitude of freezing point depression?
\item Why does the identity (not just the amount) of the solute matter in colligative properties?
\item How does the boiling point of a solution compare to that of the pure solvent?
\item Why does adding solute increase the boiling point of a liquid?
\item What physical property remains unchanged when a solution’s freezing point is lowered?
\item Why does temperature itself not change the amount of freezing point depression?
\item Why do electrolytes and nonelectrolytes differ in their effects on colligative properties?
\item How does a neutralization reaction differ from a precipitation or redox reaction?
\item Why is acid–base neutralization classified as a double replacement reaction?
\item What distinguishes a synthesis reaction from a decomposition reaction?
\item How is a combustion reaction identified based on its reactants and products?
\item Why does every combustion reaction of a hydrocarbon produce CO$_2$ and H$_2$O?
\item How does a single replacement reaction differ from a double replacement reaction?
\item Why must the activity of elements be considered when predicting single replacement outcomes?
\item How can you identify the limiting reactant in a reaction from given amounts of reactants?
\item Why does conservation of mass require balancing chemical equations?
\item How can stoichiometric ratios be used to determine theoretical yield?
\item Why do ionic compounds dissociate in water while molecular ones may not?
\item How does solubility influence whether a precipitate forms in a double replacement reaction?
\item Why is water often called the “universal solvent” in ionic and polar systems?
\item How does polarity explain why oil and water do not mix?
\item Why do stronger intermolecular forces lead to higher boiling points?
\item How can temperature changes indicate whether a process is endothermic or exothermic?
\item Why is combustion considered an exothermic oxidation process?
\item What evidence suggests that a chemical reaction, not a physical change, has occurred?
\end{itemize}\\
\hline
\end{tabular}


\newpage

\section*{\underline{Activity 1}}
\vspace{0.2cm}
\begin{tabular}{|p{4.2cm}|p{9.8cm}|}
\hline
\textbf{Collaborative Learning Techniques} & Jigsaw\\
\hline
\textbf{Learning Strategy} & Predict Test Questions \\
\hline
\textbf{Time} & 12:10---12:30\\
\hline
\textbf{Materials} & Pencil, eraser, calculator, periodic table \\
\hline
\textbf{Lecture/Textbook/Old Plans/Slide References} & Chapter 8: Solutions, Chapter 9: Chemical Reactions \\
\hline
\end{tabular}
\vspace{0.5cm}

{\large\textbf{Definition}}\\
\noindent
\begin{tabular}{|p{0.9\linewidth}|}
\hline
\begin{enumerate}
    \item\textbf{Jigsaw:} Jigsaws, when used properly, make the group as a whole dependent upon all of them in subgroups. Each group provides a peice of the puzzle. Group members are broken into smaller groups. Each small group works on some aspect of the same problem, question, or issue. They then share their part of the puzzle with the large group.
    \item\textbf{Predict Test Questions:}  Put students in groups of 2 or 3 and assign them to write a test question for a specific topic, ensuring that all topics have been convered. Ask students to write their questions on the board or on an overhead for discussion (would the professor ask the question? What is the answer? etc.). Students will have the benefit of learning to think like the teacher and they'll be able to see additional questions that other students have writen.
 
\end{enumerate} \\
\hline
\end{tabular}
\vspace{0.2cm}
\newpage
{\large\textbf{Instructions (please provide a\hl{DETAILED} activity breakdown below (and/or attached)}}\\
\noindent
\begin{tabular}{|p{0.9\linewidth}|}
\hline
\vspace{0.2cm}
\begin{enumerate}
    \item \textbf{Distribute Worksheets:} Each student receives a worksheet containing chemistry questions. The problems cover a range of difficulty so that students can decide which ones best match their needs.
    
    \item \textbf{Student Choice of Starting Point:} Students look over the worksheet and select their own starting problem rather than being assigned a fixed order. This allows them to challenge themselves where they feel it is most beneficial.
    
    \item \textbf{Staggered Timing:} Students work in pairs. Each partner begins at a different point (for example, one may start at Problem 1 while the other starts at Problem 3). They continue by alternating every three questions. This staggered structure ensures that students will be asking for help or requesting solutions at different times, keeping the class flow balanced.
    
    \item \textbf{Independent Work:} Students attempt their chosen problems individually, writing down both reasoning and answers. Emphasis is placed on documenting thought process as well as the result.
    
    \item \textbf{Help and Solution Reveal:} After the set time for each problem expires, the SI leader either walks through the solution on the whiteboard or directs students to check the answer key. Students are encouraged to ask clarifying questions at the timed checkpoints rather than all at once.
    
    \item \textbf{Partner Collaboration:} Once a solution is reviewed, partners share their reasoning for their respective problems. They ask each other questions, offer corrections, and compare how their approaches aligned with the solution.
    
    \item \textbf{Rotation of Problems:} After every three questions, students switch to new problems, continuing the staggered order so that they stay offset from each other. This ensures that each student engages with a variety of problems across the worksheet.
    
    \item \textbf{Class-Wide Debrief:} The SI leader pauses periodically to highlight frequently asked questions, difficult concepts, or common errors observed during the staggered work. Whole-class solutions are discussed when many students identify the same challenge.
    
    \item \textbf{Reflection:} At the close, students note which problems gave them the most insight, what strategies helped, and what topics they still want clarified. This reflection helps focus future review sessions.
\end{enumerate} \\
\hline
\end{tabular}
\newpage

\noindent
{\large\textbf{Content (fully worked out problems \& answers) continued...}}\\[0.2em]
\vspace{0.2cm}



\noindent\begin{tabular}{|p{0.9\linewidth}|}
\hline
\textbf{(1)} Which of the following changes how much a solution’s freezing point is lowered?\\[0.2cm]
\begin{minipage}[t]{0.85\linewidth}
\begin{enumerate}[label=\alph*.]
\item How much solute is added
\item The type of solute (number of particles it makes)
\item A \& B
\item The color of the solute
\end{enumerate}
\end{minipage}\\[0.2cm]
\hline
\end{tabular}

\noindent\begin{tabular}{|p{0.9\linewidth}|}
\hline
\textbf{(2)} What type of reaction occurs when methane burns in oxygen to produce carbon dioxide and water?\\[0.2cm]
CH$_4$ + 2O$_2$ $\rightarrow$ CO$_2$ + 2H$_2$O\\[0.2cm]
\begin{minipage}[t]{0.85\linewidth}
\begin{enumerate}[label=\alph*.]
\item Combination
\item Decomposition
\item Combustion
\item Single displacement
\end{enumerate}
\end{minipage}\\[0.2cm]
\hline
\end{tabular}

\noindent\begin{tabular}{|p{0.9\linewidth}|}
\hline
\textbf{(3)} What type of reaction occurs when potassium iodide reacts with lead(II) nitrate to form lead(II) iodide and potassium nitrate?\\[0.2cm]
2KI + Pb(NO$_3$)$_2$ $\rightarrow$ 2KNO$_3$ + PbI$_2$\\[0.2cm]
\begin{minipage}[t]{0.85\linewidth}
\begin{enumerate}[label=\alph*.]
\item Exchange
\item Combination
\item Decomposition
\item Combustion
\end{enumerate}
\end{minipage}\\[0.2cm]
\hline
\end{tabular}

\noindent\begin{tabular}{|p{0.9\linewidth}|}
\hline
\textbf{(4)} Vapor pressure lowering is an example of what type of property?\\[0.2cm]
\begin{minipage}[t]{0.85\linewidth}
\begin{enumerate}[label=\alph*.]
\item Colligative property
\item Chemical property
\item Physical property unrelated to concentration
\item Intensive property
\end{enumerate}
\end{minipage}\\[0.2cm]
\hline
\end{tabular}
\newpage

\noindent\begin{tabular}{|p{0.9\linewidth}|}
\hline
\textbf{(5)} Identify the reaction type:\\[0.2cm]
2H$_2$ + O$_2$ $\rightarrow$ 2H$_2$O\\[0.2cm]
\begin{minipage}[t]{0.85\linewidth}
\begin{enumerate}[label=\alph*.]
\item Exchage
\item Combination
\item Combustion
\item Decomposition
\end{enumerate}
\end{minipage}\\[0.2cm]
\hline
\end{tabular}

\noindent\begin{tabular}{|p{0.9\linewidth}|}
\hline
\textbf{(6)} In a redox reaction, the oxidizing agent is the species that:\\[0.2cm]
\begin{minipage}[t]{0.85\linewidth}
\begin{enumerate}[label=\alph*.]
\item Loses electrons
\item Gains electrons
\item Is always neutral
\item Has the highest atomic mass
\end{enumerate}
\end{minipage}\\[0.2cm]
\hline
\end{tabular}

\noindent\begin{tabular}{|p{0.9\linewidth}|}
\hline
\textbf{(7)} What type of chemical reaction occurs between sodium phosphate and calcium nitrate to form calcium phosphate and sodium nitrate?\\[0.2cm]
2Na$_3$PO$_4$ + 3Ca(NO$_3$)$_2$ $\rightarrow$ Ca$_3$(PO$_4$)$_2$ + 6NaNO$_3$\\[0.2cm]
\begin{minipage}[t]{0.85\linewidth}
\begin{enumerate}[label=\alph*.]
\item Decomposition
\item Exchange
\item Combination
\item Combustion
\end{enumerate}
\end{minipage}\\[0.2cm]
\hline
\end{tabular}

\noindent\begin{tabular}{|p{0.9\linewidth}|}
\hline
\textbf{(8)} Which of the following would have the highest osmolarity in water?\\[0.2cm]
\begin{minipage}[t]{0.85\linewidth}
\begin{enumerate}[label=\alph*.]
\item 1~M NaCl
\item 1~M glucose
\item 1~M MgCl$_2$
\item 1~M urea
\end{enumerate}
\end{minipage}\\[0.2cm]
\hline
\end{tabular}
\newpage


\noindent\begin{tabular}{|p{0.9\linewidth}|}
\hline
\textbf{(9)} A substance that gains electrons during a reaction is:\\[0.2cm]
\begin{minipage}[t]{0.85\linewidth}
\begin{enumerate}[label=\alph*.]
\item Oxidized
\item Reduced
\item Neutralized
\item Dissociated
\end{enumerate}
\end{minipage}\\[0.2cm]
\hline
\end{tabular}

\noindent\begin{tabular}{|p{0.9\linewidth}|}
\hline
\textbf{(10)} If two solutions have the same osmotic pressure, they are said to be:\\[0.2cm]
\begin{minipage}[t]{0.85\linewidth}
\begin{enumerate}[label=\alph*.]
\item Hypertonic
\item Hypotonic
\item Isotonic
\item Colloidal
\end{enumerate}
\end{minipage}\\[0.2cm]
\hline
\end{tabular}

\noindent\begin{tabular}{|p{0.9\linewidth}|}
\hline
\textbf{(11)} What type of reaction takes place when sodium carbonate reacts with calcium chloride to form calcium carbonate and sodium chloride?\\[0.2cm]
Na$_2$CO$_3$ + CaCl$_2$ $\rightarrow$ CaCO$_3$ + 2NaCl\\[0.2cm]
\begin{minipage}[t]{0.85\linewidth}
\begin{enumerate}[label=\alph*.]
\item Combination
\item Exchange
\item Single replacement
\item Decomposition
\end{enumerate}
\end{minipage}\\[0.2cm]
\hline
\end{tabular}

\noindent\begin{tabular}{|p{0.9\linewidth}|}
\hline
\textbf{(12)} In a redox reaction, the reducing agent is the substance that:\\[0.2cm]
\begin{minipage}[t]{0.85\linewidth}
\begin{enumerate}[label=\alph*.]
\item Gains electrons
\item Loses electrons
\item Is neither oxidized nor reduced
\item Lowers its oxidation number
\end{enumerate}
\end{minipage}\\[0.2cm]
\hline
\end{tabular}
\newpage


\noindent\begin{tabular}{|p{0.9\linewidth}|}
\hline
\textbf{(13)} What type of chemical reaction does ethanol undergo in air?\\[0.2cm]
C$_2$H$_5$OH + 3O$_2$ $\rightarrow$ 2CO$_2$ + 3H$_2$O\\[0.2cm]
\begin{minipage}[t]{0.85\linewidth}
\begin{enumerate}[label=\alph*.]
\item Neutralization
\item Single replacement
\item Combustion
\item Precipitation
\end{enumerate}
\end{minipage}\\[0.2cm]
\hline
\end{tabular}

\noindent\begin{tabular}{|p{0.9\linewidth}|}
\hline
\textbf{(14)} When a solute is added to a solvent, what happens to the freezing point of the solution?\\[0.2cm]
\begin{minipage}[t]{0.85\linewidth}
\begin{enumerate}[label=\alph*.]
\item It increases
\item It decreases
\item It stays the same
\item It becomes zero
\end{enumerate}
\end{minipage}\\[0.2cm]
\hline
\end{tabular}

\noindent\begin{tabular}{|p{0.9\linewidth}|}
\hline
\textbf{(15)} In the reaction Zn + Cu$^{2+}$ $\rightarrow$ Zn$^{2+}$ + Cu, which species is oxidized?\\[0.2cm]
\begin{minipage}[t]{0.85\linewidth}
\begin{enumerate}[label=\alph*.]
\item Cu$^{2+}$
\item Cu
\item Zn
\item H$_2$O
\end{enumerate}
\end{minipage}\\[0.2cm]
\hline
\end{tabular}

\noindent\begin{tabular}{|p{0.9\linewidth}|}
\hline
\textbf{(16)} Which factor determines how much the freezing point of a solution is lowered?\\[0.2cm]
\begin{minipage}[t]{0.85\linewidth}
\begin{enumerate}[label=\alph*.]
\item The number of solute particles
\item The solute's color
\item The solvent's taste
\item The solute's shape
\end{enumerate}
\end{minipage}\\[0.2cm]
\hline
\end{tabular}
\newpage


\noindent\begin{tabular}{|p{0.9\linewidth}|}
\hline
\textbf{(17)} What type of reaction occurs when acetylene reacts with oxygen to produce CO$_2$ and H$_2$O?\\[0.2cm]
2C$_2$H$_2$ + 5O$_2$ $\rightarrow$ 4CO$_2$ + 2H$_2$O\\[0.2cm]
\begin{minipage}[t]{0.85\linewidth}
\begin{enumerate}[label=\alph*.]
\item Neutralization
\item Combustion
\item Combination
\item Decomposition
\end{enumerate}
\end{minipage}\\[0.2cm]
\hline
\end{tabular}

\noindent\begin{tabular}{|p{0.9\linewidth}|}
\hline
\textbf{(18)} Adding salt to water causes which of the following effects?\\[0.2cm]
\begin{minipage}[t]{0.85\linewidth}
\begin{enumerate}[label=\alph*.]
\item Raises freezing point and lowers boiling point
\item Lowers freezing point and raises boiling point
\item Raises both freezing and boiling points
\item Lowers both freezing and boiling points
\end{enumerate}
\end{minipage}\\[0.2cm]
\hline
\end{tabular}

\noindent\begin{tabular}{|p{0.9\linewidth}|}
\hline
\textbf{(19)} What happens to the vapor pressure of a solvent when a nonvolatile solute is added?\\[0.2cm]
\begin{minipage}[t]{0.85\linewidth}
\begin{enumerate}[label=\alph*.]
\item It increases
\item It decreases
\item It stays the same
\item It becomes equal to the boiling point
\end{enumerate}
\end{minipage}\\[0.2cm]
\hline
\end{tabular}

\noindent\begin{tabular}{|p{0.9\linewidth}|}
\hline
\textbf{(20)} Identify the type of reaction below:\\[0.2cm]
K$_2$SO$_4$ + Ba(NO$_3$)$_2$ $\rightarrow$ 2KNO$_3$ + BaSO$_4$\\[0.2cm]
\begin{minipage}[t]{0.85\linewidth}
\begin{enumerate}[label=\alph*.]
\item Exchange
\item Combination
\item Combustion
\item Decomposition
\end{enumerate}
\end{minipage}\\[0.2cm]
\hline
\end{tabular}
\newpage


\noindent\begin{tabular}{|p{0.9\linewidth}|}
\hline
\textbf{(21)} Which solution would have the lowest vapor pressure?\\[0.2cm]
\begin{minipage}[t]{0.85\linewidth}
\begin{enumerate}[label=\alph*.]
\item Pure water
\item 1~M glucose solution
\item 1~M NaCl solution
\item 0.5~M NaCl solution
\end{enumerate}
\end{minipage}\\[0.2cm]
\hline
\end{tabular}

\noindent\begin{tabular}{|p{0.9\linewidth}|}
\hline
\textbf{(22)} Which of the following elements always has an oxidation number of +1 in compounds?\\[0.2cm]
\begin{minipage}[t]{0.85\linewidth}
\begin{enumerate}[label=\alph*.]
\item Sodium
\item Oxygen
\item Fluorine
\item Sulfur
\end{enumerate}
\end{minipage}\\[0.2cm]
\hline
\end{tabular}

\noindent\begin{tabular}{|p{0.9\linewidth}|}
\hline
\textbf{(23)} Identify the reaction type:\\[0.2cm]
2H$_2$O$_2$ $\rightarrow$ 2H$_2$O + O$_2$\\[0.2cm]
\begin{minipage}[t]{0.85\linewidth}
\begin{enumerate}[label=\alph*.]
\item Combination
\item Combustion
\item Decomposition
\item Exchange
\end{enumerate}
\end{minipage}\\[0.2cm]
\hline
\end{tabular}


\noindent\begin{tabular}{|p{0.9\linewidth}|}
\hline
\textbf{(24)} When hydrochloric acid reacts with sodium hydroxide to produce sodium chloride and water, what kind of reaction is it?\\[0.2cm]
HCl + NaOH $\rightarrow$  NaCl + H$_2$O\\[0.2cm]
\begin{minipage}[t]{0.85\linewidth}
\begin{enumerate}[label=\alph*.]
\item Combustion
\item Exchange
\item Combination
\item Decomposition
\end{enumerate}
\end{minipage}\\[0.2cm]
\hline
\end{tabular}
\newpage


\noindent\begin{tabular}{|p{0.9\linewidth}|}
\hline
\textbf{(25)} A 1.0~m solution of NaCl has a higher boiling point than a 1.0~m solution of glucose because:\\[0.2cm]
\begin{minipage}[t]{0.85\linewidth}
\begin{enumerate}[label=\alph*.]
\item NaCl has a smaller molar mass
\item NaCl dissociates into more particles
\item Glucose has more hydrogen bonds
\item Glucose is more polar
\end{enumerate}
\end{minipage}\\[0.2cm]
\hline
\end{tabular}

\noindent\begin{tabular}{|p{0.9\linewidth}|}
\hline
\textbf{(26)} What term describes the elevation of a solvent's boiling point when solute is added?\\[0.2cm]
\begin{minipage}[t]{0.85\linewidth}
\begin{enumerate}[label=\alph*.]
\item Freezing point depression
\item Boiling point elevation
\item Osmotic pressure
\item Vapor pressure lowering
\end{enumerate}
\end{minipage}\\[0.2cm]
\hline
\end{tabular}

\noindent\begin{tabular}{|p{0.9\linewidth}|}
\hline
\textbf{(27)} What kind of reaction is the burning of glucose in oxygen to form carbon dioxide and water?\\[0.2cm]
C$_6$H$_{12}$O$_6$ + 6O$_2$ $\rightarrow$  6CO$_2$ + 6H$_2$O\\[0.2cm]
\begin{minipage}[t]{0.85\linewidth}
\begin{enumerate}[label=\alph*.]
\item Double displacement
\item Combustion
\item Decomposition
\item Combination
\end{enumerate}
\end{minipage}\\[0.2cm]
\hline
\end{tabular}

\noindent\begin{tabular}{|p{0.9\linewidth}|}
\hline
\textbf{(28)} Which element is being reduced in the reaction Fe$_2$O$_3$ + 3CO $\rightarrow$ 2Fe + 3CO$_2$?\\[0.2cm]
\begin{minipage}[t]{0.85\linewidth}
\begin{enumerate}[label=\alph*.]
\item Fe
\item O
\item C
\item CO$_2$
\end{enumerate}
\end{minipage}\\[0.2cm]
\hline
\end{tabular}
\newpage


\noindent\begin{tabular}{|p{0.9\linewidth}|}
\hline
\textbf{(29)} When ammonium chloride reacts with silver nitrate to form silver chloride and ammonium nitrate, the reaction is classified as:\\[0.2cm]
NH$_4$Cl + AgNO$_3$ $\rightarrow$ AgCl + NH$_4$NO$_3$\\[0.2cm]
\begin{minipage}[t]{0.85\linewidth}
\begin{enumerate}[label=\alph*.]
\item Combination
\item Exchange
\item Decomposition
\item Combustion
\end{enumerate}
\end{minipage}\\[0.2cm]
\hline
\end{tabular}

\noindent\begin{tabular}{|p{0.9\linewidth}|}
\hline
\textbf{(30)} The decrease in vapor pressure when solute is added is known as:\\[0.2cm]
\begin{minipage}[t]{0.85\linewidth}
\begin{enumerate}[label=\alph*.]
\item Vapor pressure elevation
\item Vapor pressure lowering
\item Freezing point depression
\item Boiling point elevation
\end{enumerate}
\end{minipage}\\[0.2cm]
\hline
\end{tabular}

\noindent\begin{tabular}{|p{0.9\linewidth}|}
\hline
\textbf{(31)} When more solute is added to a solution, the osmotic pressure will:\\[0.2cm]
\begin{minipage}[t]{0.85\linewidth}
\begin{enumerate}[label=\alph*.]
\item Decrease
\item Stay the same
\item Increase
\item Become zero
\end{enumerate}
\end{minipage}\\[0.2cm]
\hline
\end{tabular}

\noindent\begin{tabular}{|p{0.9\linewidth}|}
\hline
\textbf{(32)} Osmotic pressure depends on which of the following factors?\\[0.2cm]
\begin{minipage}[t]{0.85\linewidth}
\begin{enumerate}[label=\alph*.]
\item The number of solute particles in solution
\item The color of the solute
\item The shape of the solvent molecules
\item The size of the container
\end{enumerate}
\end{minipage}\\[0.2cm]
\hline
\end{tabular}
\newpage


\noindent\begin{tabular}{|p{0.9\linewidth}|}
\hline
\textbf{(33)} What is the oxidation number of chlorine in NaCl?\\[0.2cm]
\begin{minipage}[t]{0.85\linewidth}
\begin{enumerate}[label=\alph*.]
\item +1
\item -1
\item 0
\item +2
\end{enumerate}
\end{minipage}\\[0.2cm]
\hline
\end{tabular}

\noindent\begin{tabular}{|p{0.9\linewidth}|}
\hline
\textbf{(34)} What type of reaction occurs when silver nitrate reacts with sodium chloride to produce silver chloride and sodium nitrate?\\[0.2cm]
AgNO$_3$ + NaCl $\rightarrow$ AgCl + NaNO$_3$\\[0.2cm]
\begin{minipage}[t]{0.85\linewidth}
\begin{enumerate}[label=\alph*.]
\item Combination
\item Decomposition
\item Exchange
\item Combustion
\end{enumerate}
\end{minipage}\\[0.2cm]
\hline
\end{tabular}

\noindent\begin{tabular}{|p{0.9\linewidth}|}
\hline
\textbf{(35)} What happens to the boiling point of a solvent when a nonvolatile solute is added?\\[0.2cm]
\begin{minipage}[t]{0.85\linewidth}
\begin{enumerate}[label=\alph*.]
\item It decreases
\item It remains the same
\item It increases
\item It fluctuates randomly
\end{enumerate}
\end{minipage}\\[0.2cm]
\hline
\end{tabular}

\noindent\begin{tabular}{|p{0.9\linewidth}|}
\hline
\textbf{(36)} When sodium reacts with chlorine to produce sodium chloride, what type of reaction is this?\\[0.2cm]
2Na + Cl$_2$ $\rightarrow$ 2NaCl\\[0.2cm]
\begin{minipage}[t]{0.85\linewidth}
\begin{enumerate}[label=\alph*.]
\item Decomposition
\item Combustion
\item Exchange
\item Combination
\end{enumerate}
\end{minipage}\\[0.2cm]
\hline
\end{tabular}
\newpage

\noindent\begin{tabular}{|p{0.9\linewidth}|}
\hline
\textbf{(37)} What happens to an atom's oxidation number when it is oxidized?\\[0.2cm]
\begin{minipage}[t]{0.85\linewidth}
\begin{enumerate}[label=\alph*.]
\item It increases
\item It decreases
\item It stays the same
\item It becomes zero
\end{enumerate}
\end{minipage}\\[0.2cm]
\hline
\end{tabular}

\noindent\begin{tabular}{|p{0.9\linewidth}|}
\hline
\textbf{(38)} What type of reaction occurs when lithium sulfate reacts with barium nitrate to form barium sulfate and lithium nitrate?\\[0.2cm]
Li$_2$SO$_4$ + Ba(NO$_3$)$_2$ $\rightarrow$ BaSO$_4$ + 2LiNO$_3$\\[0.2cm]
\begin{minipage}[t]{0.85\linewidth}
\begin{enumerate}[label=\alph*.]
\item Combination
\item Decomposition
\item Exchange
\item Combustion
\end{enumerate}
\end{minipage}\\[0.2cm]
\hline
\end{tabular}

\noindent\begin{tabular}{|p{0.9\linewidth}|}
\hline
\textbf{(39)} What is the oxidation number of nitrogen in NH$_3$?\\[0.2cm]
\begin{minipage}[t]{0.85\linewidth}
\begin{enumerate}[label=\alph*.]
\item +3
\item -3
\item 0
\item +5
\end{enumerate}
\end{minipage}\\[0.2cm]
\hline
\end{tabular}

\noindent\begin{tabular}{|p{0.9\linewidth}|}
\hline
\textbf{(40)} Identify the reaction type:\\[0.2cm]
BaCl$_2$ + Na$_2$SO$_4$ $\rightarrow$ BaSO$_4$ + 2NaCl\\[0.2cm]
\begin{minipage}[t]{0.85\linewidth}
\begin{enumerate}[label=\alph*.]
\item Combination
\item Decomposition
\item Exchange
\item Single replacement
\end{enumerate}
\end{minipage}\\[0.2cm]
\hline
\end{tabular}
\newpage


\noindent\begin{tabular}{|p{0.9\linewidth}|}
\hline
\textbf{(41)} What happens to osmotic pressure when temperature increases (all else equal)?\\[0.2cm]
\begin{minipage}[t]{0.85\linewidth}
\begin{enumerate}[label=\alph*.]
\item It decreases
\item It increases
\item It stays constant
\item It becomes zero
\end{enumerate}
\end{minipage}\\[0.2cm]
\hline
\end{tabular}

\noindent\begin{tabular}{|p{0.9\linewidth}|}
\hline
\textbf{(42)} The burning of benzene in air to form carbon dioxide and water is an example of what type of reaction?\\[0.2cm]
2C$_6$H$_6$ + 15O$_2$ $\rightarrow$ 12CO$_2$ + 6H$_2$O\\[0.2cm]
\begin{minipage}[t]{0.85\linewidth}
\begin{enumerate}[label=\alph*.]
\item Oxidation-reduction
\item Combustion
\item Decomposition
\item Neutralization
\end{enumerate}
\end{minipage}\\[0.2cm]
\hline
\end{tabular}

\noindent\begin{tabular}{|p{0.9\linewidth}|}
\hline
\textbf{(43)} The depression of the freezing point is directly proportional to which quantity?\\[0.2cm]
\begin{minipage}[t]{0.85\linewidth}
\begin{enumerate}[label=\alph*.]
\item Solute's molar mass
\item Solvent's density
\item Number of solute particles in solution
\item Volume of solvent
\end{enumerate}
\end{minipage}\\[0.2cm]
\hline
\end{tabular}

\noindent\begin{tabular}{|p{0.9\linewidth}|}
\hline
\textbf{(44)} Which factor mainly determines the amount of vapor pressure lowering in a solution?\\[0.2cm]
\begin{minipage}[t]{0.85\linewidth}
\begin{enumerate}[label=\alph*.]
\item The number of solute particles
\item The color of the solute
\item The shape of the solvent molecules
\item The pressure of the air above the solution
\end{enumerate}
\end{minipage}\\[0.2cm]
\hline
\end{tabular}
\newpage

\noindent\begin{tabular}{|p{0.9\linewidth}|}
\hline
\textbf{(45)} In the reaction 2Na + Cl$_2$ $\rightarrow$ 2NaCl, which is the oxidizing agent?\\[0.2cm]
\begin{minipage}[t]{0.85\linewidth}
\begin{enumerate}[label=\alph*.]
\item Na
\item Cl$_2$
\item NaCl
\item Both Na and Cl$_2$
\end{enumerate}
\end{minipage}\\[0.2cm]
\hline
\end{tabular}

\noindent\begin{tabular}{|p{0.9\linewidth}|}
\hline
\textbf{(46)} Osmosis is the movement of solvent molecules from:\\[0.2cm]
\begin{minipage}[t]{0.85\linewidth}
\begin{enumerate}[label=\alph*.]
\item Low solute concentration to high solute concentration
\item High solute concentration to low solute concentration
\item High pressure to low pressure
\item High temperature to low temperature
\end{enumerate}
\end{minipage}\\[0.2cm]
\hline
\end{tabular}

\noindent\begin{tabular}{|p{0.9\linewidth}|}
\hline
\textbf{(47)} What is osmotic pressure?\\[0.2cm]
\begin{minipage}[t]{0.85\linewidth}
\begin{enumerate}[label=\alph*.]
\item The pressure needed to stop osmosis
\item The pressure caused by boiling
\item The force between solute and solvent
\item The vapor pressure of a liquid
\end{enumerate}
\end{minipage}\\[0.2cm]
\hline
\end{tabular}

\noindent\begin{tabular}{|p{0.9\linewidth}|}
\hline
\textbf{(48)} Which of the following reactions is a redox reaction?\\[0.2cm]
\begin{minipage}[t]{0.85\linewidth}
\begin{enumerate}[label=\alph*.]
\item Zn + CuSO$_4$ $\rightarrow$ ZnSO$_4$ + Cu\\[0.1cm]
\item AgNO$_3$ + NaCl $\rightarrow$ AgCl + NaNO$_3$\\[0.1cm]
\item CaCO$_3$ $\rightarrow$ CaO + CO$_2$\\[0.1cm]
\item H$_2$O + SO$_3$ $\rightarrow$ H$_2$SO$_4$\\[0.1cm]
\end{enumerate}
\end{minipage}\\[0.2cm]
\hline
\end{tabular}
\newpage

\noindent\begin{tabular}{|p{0.9\linewidth}|}
\hline
\textbf{(49)} What happens to a cell placed in a hypotonic solution?\\[0.2cm]
\begin{minipage}[t]{0.85\linewidth}
\begin{enumerate}[label=\alph*.]
\item It shrinks
\item It swells
\item It stays the same size
\item It loses solute
\end{enumerate}
\end{minipage}\\[0.2cm]
\hline
\end{tabular}

\noindent\begin{tabular}{|p{0.9\linewidth}|}
\hline
\textbf{(50)} Which of the following increases osmotic pressure?\\[0.2cm]
\begin{minipage}[t]{0.85\linewidth}
\begin{enumerate}[label=\alph*.]
\item Lower solute concentration
\item Higher solute concentration
\item Larger solvent molecules
\item Higher freezing point
\end{enumerate}
\end{minipage}\\[0.2cm]
\hline
\end{tabular}

\noindent\begin{tabular}{|p{0.9\linewidth}|}
\hline
\textbf{(51)} Why does adding a solute lower the vapor pressure of a solvent?\\[0.2cm]
\begin{minipage}[t]{0.85\linewidth}
\begin{enumerate}[label=\alph*.]
\item Fewer solvent molecules can escape into the gas phase
\item Solute molecules turn into vapor more easily
\item The solute increases the temperature of the liquid
\item The solvent molecules become heavier
\end{enumerate}
\end{minipage}\\[0.2cm]
\hline
\end{tabular}

\noindent\begin{tabular}{|p{0.9\linewidth}|}
\hline
\textbf{(52)} What is the oxidation number of manganese in KMnO$_4$?\\[0.2cm]
\begin{minipage}[t]{0.85\linewidth}
\begin{enumerate}[label=\alph*.]
\item +2
\item +4
\item +6
\item +7
\end{enumerate}
\end{minipage}\\[0.2cm]
\hline
\end{tabular}
\newpage

\noindent\begin{tabular}{|p{0.9\linewidth}|}
\hline
\textbf{(53)} What is the oxidation number of sulfur in SO$_2$?\\[0.2cm]
\begin{minipage}[t]{0.85\linewidth}
\begin{enumerate}[label=\alph*.]
\item +2
\item +4
\item +6
\item 0
\end{enumerate}
\end{minipage}\\[0.2cm]
\hline
\end{tabular}

\noindent\begin{tabular}{|p{0.9\linewidth}|}
\hline
\textbf{(54)} If 0.5~mol of NaCl is dissolved in 1.0~kg of water, how many moles of particles are present assuming complete dissociation?\\[0.2cm]
\begin{minipage}[t]{0.85\linewidth}
\begin{enumerate}[label=\alph*.]
\item 0.5
\item 1.0
\item 1.5
\item 2.0
\end{enumerate}
\end{minipage}\\[0.2cm]
\hline
\end{tabular}

\noindent\begin{tabular}{|p{0.9\linewidth}|}
\hline
\textbf{(55)} The addition of a solute to a solvent causes the boiling point to:\\[0.2cm]
\begin{minipage}[t]{0.85\linewidth}
\begin{enumerate}[label=\alph*.]
\item Decrease
\item Increase
\item Remain constant
\item Become equal to the freezing point
\end{enumerate}
\end{minipage}\\[0.2cm]
\hline
\end{tabular}

\noindent\begin{tabular}{|p{0.9\linewidth}|}
\hline
\textbf{(56)} Which of the following represents a redox reaction?\\[0.2cm]
\begin{minipage}[t]{0.85\linewidth}
\begin{enumerate}[label=\alph*.]
\item HCl + NaOH $\rightarrow$ NaCl + H$_2$O\\[0.1cm]
\item Zn + 2HCl $\rightarrow$ ZnCl$_2$ + H$_2$\\[0.1cm]
\item CaCO$_3$ $\rightarrow$ CaO + CO$_2$\\[0.1cm]
\item NaOH + HNO$_3$ $\rightarrow$ NaNO$_3$ + H$_2$O\\[0.1cm]
\end{enumerate}

\end{minipage}\\[0.2cm]
\hline
\end{tabular}
\newpage

\noindent\begin{tabular}{|p{0.9\linewidth}|}
\hline
\textbf{(57)} Which of the following would cause the largest increase in the boiling point of water?\\[0.2cm]
\begin{minipage}[t]{0.85\linewidth}
\begin{enumerate}[label=\alph*.]
\item 1~mol NaCl
\item 1~mol glucose
\item 1~mol MgCl$_2$
\item 1~mol urea
\end{enumerate}
\end{minipage}\\[0.2cm]
\hline
\end{tabular}

\noindent\begin{tabular}{|p{0.9\linewidth}|}
\hline
\textbf{(58)} What is the oxidation number of oxygen in H$_2$O?\\[0.2cm]
\begin{minipage}[t]{0.85\linewidth}
\begin{enumerate}[label=\alph*.]
\item 0
\item -1
\item -2
\item +2
\end{enumerate}
\end{minipage}\\[0.2cm]
\hline
\end{tabular}

\noindent\begin{tabular}{|p{0.9\linewidth}|}
\hline
\textbf{(59)} The reaction between calcium hydroxide and hydrochloric acid produces calcium chloride and water. Identify the reaction type.\\[0.2cm]
Ca(OH)$_2$ + 2HCl $\rightarrow$ CaCl$_2$ + 2H$_2$O\\[0.2cm]
\begin{minipage}[t]{0.85\linewidth}
\begin{enumerate}[label=\alph*.]
\item Decomposition
\item Exchange
\item Combustion
\item Combination
\end{enumerate}
\end{minipage}\\[0.2cm]
\hline
\end{tabular}

\noindent\begin{tabular}{|p{0.9\linewidth}|}
\hline
\textbf{(60)} Osmolarity refers to:\\[0.2cm]
\begin{minipage}[t]{0.85\linewidth}
\begin{enumerate}[label=\alph*.]
\item The number of solute particles per liter of solution
\item The mass of solute per liter of solvent
\item The number of moles per kilogram of solvent
\item The ratio of solvent to solute molecules
\end{enumerate}
\end{minipage}\\[0.2cm]
\hline
\end{tabular}

\newpage
{\large\textbf{Checking for Understanding (general conceptual questions)}}\\
\noindent \begin{tabular}{|p{0.9\linewidth}|}
\hline
\vspace{0.2cm}

\begin{enumerate}
    \item \textbf{Targeted focus through self-selection:} When students choose which questions to tackle from the worksheet, they naturally gravitate toward areas they feel less confident about, ensuring the group addresses gaps in understanding.
    \item \textbf{Peer modeling of problem-solving:} Seeing a peer work through a problem on the board exposes students to different ways of setting up and explaining a solution, clarifying concepts for observers.
    \item \textbf{Reflection on understanding:} Students must decide whether they can solve a question themselves or need help, prompting metacognition and highlighting what they truly understand versus what requires clarification.
    \item \textbf{Immediate correction of misconceptions:} Peers correcting or expanding an explanation in real time prevents misunderstandings from becoming ingrained.
    \item \textbf{Balancing independence and collaboration:} Students first work independently on a question, then collaborate during discussion, ensuring individual accountability while leveraging group problem-solving.
    \item \textbf{Distinguishing conceptual vs. procedural questions:} Students can classify questions by identifying key ideas versus stepwise calculations, determining whether a question emphasizes big-picture reasoning or step-by-step execution.
    \item \textbf{Connecting big concepts to details:} Rotating through different student-selected questions helps the group link overarching concepts to multiple specific examples, reinforcing understanding of the framework behind the problems.
    \item \textbf{Verbalizing reasoning:} Explaining steps out loud forces students to articulate assumptions, logic, and chemical principles, improving clarity and communication of ideas.
    \item \textbf{Learning from multiple approaches:} Different students may solve the same problem uniquely, exposing peers to alternative methods or perspectives and broadening their problem-solving toolkit.
    \item \textbf{Application of theory to practice:} Working collaboratively on worksheet problems allows students to connect abstract chemical principles to concrete calculations or predictions.
    \item \textbf{Instructor insight:} Questions students select or ask reveal where common difficulties lie, helping instructors or SI leaders provide targeted guidance.
    \item \textbf{Strengthening long-term understanding:} Actively solving, explaining, and correcting questions in real time consolidates memory and understanding more effectively than passive review.
\end{enumerate} \\
\hline
\end{tabular}

\newpage
\noindent
\newpage
\noindent
\section*{\underline{Activity 2}}
\begin{tabular}{|p{4.2cm}|p{9.8cm}|}
\hline
\textbf{Collaborative Learning Techniques} & Clusters \\
\hline
\textbf{Learning Strategy} & Escape Room \\
\hline
\textbf{Time} & 12:35---12:50 \\
\hline
\textbf{Materials} & A periodic table, the textbook, readiness to work as a team or share ideas \\
\hline
\textbf{Lecture/Textbook/Old Plans/Slide References} & Chapter 7: Gases, Liquids and Solids, Chapter 8: Solutions \\
\hline
\end{tabular}

\vspace{0.5cm}
\noindent
\begin{tabular}{|p{0.9\linewidth}|}
\hline
{\large\textbf{Definitions}}
\begin{enumerate}
    \item \textbf{Escape Room:} Students enter a themed, time-bound challenge where they must solve a series of puzzles, riddles, or tasks to “escape” the room. Each puzzle is designed to test understanding of specific course concepts, requiring critical thinking, teamwork, and creativity. Students can expect to encounter hidden clues, coded messages, or problem-solving stations that build on one another; progress usually depends on collaboration and discussion. The SI leader will circulate to provide subtle hints, clarify misconceptions, and ensure students stay on track without giving away answers. This activity promotes peer-to-peer learning, reinforces content in an engaging way, and encourages students to communicate their reasoning while practicing time management under pressure.
\end{enumerate} \\
\hline
\end{tabular}
\newpage

\vspace{0.5cm}
\noindent
{\large \textbf{Instructions (detailed activity breakdown below)}} \\

\begin{tabular}{|p{0.9\linewidth}|}
\hline
\begin{minipage}[t]{\linewidth}
\begin{enumerate}
    \item \textbf{Objective:} Students will collaboratively solve chemistry problems covering \textbf{molarity, dilutions, solubility rules and concentration values}. Each correct response earns points toward an overall team score. The goal is to promote discussion, problem-solving, and application of multiple concepts in a dynamic setting.

    \item \textbf{Group Size:} Pairs (2 students per team)

    \item \textbf{Setup:}
        \begin{itemize}
            \item \textbf{Station 1 – Molarity:} Explore how solute amount and solution volume determine concentration.  
            \item \textbf{Station 2 – Dilution:} Examine how adding solvent changes the molarity of a solution.  
            \item \textbf{Station 3 – Solubility:} Investigate factors that affect how much solute can dissolve in a solvent.  
            \item \textbf{Station 4 – Saturation:} Distinguish between unsaturated, saturated, and supersaturated solutions.  
            \item \textbf{Station 5 – Percent by Mass:} Analyze how to express concentration using mass of solute and total mass of solution.  
            \item \textbf{Station 6 – Mass/Volume Percent:} Study how to calculate concentration when solute is solid and solvent is liquid.  
            \item \textbf{Station 7 – Volume Percent:} Explore how to describe concentration in mixtures of liquids.  
            \item \textbf{Station 8 – Freezing Point Depression:} Examine how solutes lower the freezing point of solvents and factors affecting the magnitude.
            \item \textbf{Station 9 – Boiling Point Elevation:} Investigate how solutes increase the boiling point of solvents.  
            \item \textbf{Station 10 – Colligative Properties of Electrolytes vs. Nonelectrolytes:} Compare how ionic and molecular solutes affect freezing and boiling points differently.  
            \item \textbf{Station 11 – Reaction Types:} Identify and classify reactions (synthesis, decomposition, single replacement, double replacement, neutralization, combustion).  
            \item \textbf{Station 12 – Redox Reactions:} Identify oxidation and reduction, assign oxidation numbers, and balance redox reactions.  
            \item \textbf{Station 13 – Limiting Reactant \& Theoretical Yield:} Use stoichiometry to determine which reactant limits product formation and calculate theoretical yield.  
            \item \textbf{Station 14 – Stoichiometric Calculations:} Practice using molar ratios to convert between moles, mass, and molecules.  
            \item \textbf{Station 15 – Solvent Polarity and “Like Dissolves Like”:} Explore how solvent polarity affects solubility of different solutes.  
            \item \textbf{Station 16 – Saturated vs. Supersaturated Solutions:} Investigate conditions for forming supersaturated solutions and identifying crystallization.  
            \item \textbf{(Optional) Station 17 – Concept Review:} Summarize and connect ideas about solution concentration and solubility.
    \end{itemize}
\end{enumerate}
\end{minipage} \\
\hline
\end{tabular}
\newpage

\vspace{0.5cm}
\noindent
{\large \textbf{Instructions (detailed activity breakdown below) conitnued...}} \\

\begin{tabular}{|p{0.9\linewidth}|}
\hline
\begin{enumerate}
    \item At each station, place a set of laminated question cards (multiple choice and true/false). Each card includes:
        \begin{itemize}
            \item A unique question number and topic label.
            \item The possible answer choices (A–E or T/F).
            \item A mapping line showing which station each answer leads to (to simulate “pathways” through the escape room).
        \end{itemize}
        \item Provide each pair of students with a \textbf{team slip} (half-sheet) labeled with their team name and space to record up to 3 verified answers.
    \item \textbf{Facilitation:}
    \begin{itemize}
        \item Pairs may begin at \textbf{any station of their choice}.
        \item Each pair selects one question card, discusses the problem, and writes their chosen answer \textbf{on the back of the card}.
        \item Students then show their work and reasoning to the instructor or SI leader for quick verification.
        \item If the answer is correct, they record that problem on their team slip (problem number and answer) and may proceed to any other station of their choice.
        \item Pairs are encouraged to move freely between stations, aiming to correctly solve at least \textbf{three problems} during the activity period.
        \item After completing at least three problems, teams submit their slips to the instructor for scoring.
    \end{itemize}

    \item \textbf{Scoring:}
    \begin{itemize}
        \item Each correctly solved problem = \textbf{1 point}.
        \item Teams with the highest total score win or “escape” first.
        \item Bonus points may be awarded for clear reasoning or teamwork.
    \end{itemize}

    \item \textbf{Rules:}
    \begin{itemize}
        \item Teams must work collaboratively—no switching partners mid-activity.
        \item All work must be shown before a response is accepted.
        \item Only verified answers written on the team slip count toward the final score.
        \item Communication between pairs is encouraged, but direct sharing of answers is not permitted.
    \end{itemize}

    \item \textbf{Wrap-Up:}
    \begin{itemize}
        \item After all slips are collected, review several challenging questions as a class.
        \item Discuss common reasoning errors and strategies used by top-scoring pairs.
        \item Reinforce conceptual links between stations (e.g., gas behavior and intermolecular forces).
    \end{itemize}
\end{enumerate}
\\
\hline
\end{tabular}



\newpage
\noindent
{\large\textbf{Checking for Understanding (questions \& answers)}}\\[0.2em]
\noindent
\begin{tabular}{|p{0.9\linewidth}|}
\hline
\textbf{Questions:} \\
\begin{enumerate}
    \item \textbf{Targeted clarification:} Allowing students to bring questions from the lecture lets them focus on concepts they found most confusing.
    \item \textbf{Peer modeling of explanations:} Observing how classmates summarize or interpret lecture content helps students see alternative ways to understand the material.
    \item \textbf{Reflection on comprehension:} Discussing lecture points collaboratively encourages students to distinguish what they grasp from what needs further clarification.
    \item \textbf{Immediate correction of misconceptions:} Peers can correct or expand on explanations in real time, preventing misunderstandings from solidifying.
    \item \textbf{Balancing independent thinking and discussion:} Students first reflect individually, then contribute to group discussion, promoting both accountability and collaborative learning.
    \item \textbf{Identifying conceptual vs. procedural points:} Students can highlight which parts of the lecture are big-picture ideas versus detailed steps or examples.
    \item \textbf{Connecting concepts across topics:} Discussing multiple lecture points together helps students link overarching principles to specific details.
    \item \textbf{Verbalizing reasoning:} Explaining ideas out loud strengthens students’ ability to communicate and reason about the material clearly.
    \item \textbf{Learning from diverse perspectives:} Hearing multiple interpretations or summaries of the same lecture point exposes students to alternative approaches.
    \item \textbf{Application of theory:} Collaborative discussion encourages students to connect lecture content to examples, problems, or real-world applications.
    \item \textbf{Instructor insight:} Observing which lecture points students question most helps the instructor identify common difficulties and adjust instruction.
    \item \textbf{Reinforcing long-term understanding:} Actively reviewing and discussing lecture material together solidifies comprehension better than passive note review.
\end{enumerate} \\
\hline  
\end{tabular}

\newpage
\noindent
\section*{\underline{Closer}}
\begin{tabular}{|p{4.2cm}|p{9.8cm}|}
\hline
\textbf{Collaborative Learning Techniques} & Group Survey\\
\hline
\textbf{Learning Strategy} & Brain dump\\
\hline
\textbf{Time} & 12:50---12:55\\
\hline
\textbf{Materials} & Index cards, writing utensils\\
\hline
\textbf{Lecture/Textbook/Old Plans/Slide References} & Chapters 1-9 \\
\hline
\end{tabular}

\vspace{0.5cm}
\noindent
\begin{tabular}{|p{0.9\linewidth}|}
\hline
{\large\textbf{Definitions}}
\begin{enumerate}
   \item\textbf{Group survey:} Each group member is surveyed to discover their position on an issue, problem or topic. This process ensures that each member of the group is allowed to offer or state their point of view.
   \item\textbf{Brain dump:} This is a great closer for students to sum up their knowledge on the topic(s) from the session. Have the students take out a blank piece of paper /Google Doc and write/type("dump") everything they know relevant from the discussion that day. The SI Leader can look at their work to assess the students' understanding of the course material and have them compare it to each other. It is important to at least discuss with the students so the SI Leader knows the students' comfort levels with the material. 
\end{enumerate}
\\ \hline
\end{tabular}

\newpage
\noindent
{\large\textbf{Instructions (please provide a\hl{DETAILED} activity breakdown below (and/or attached}}

\begin{tabular}{|p{0.9\linewidth}|}
\hline
\begin{enumerate}[itemsep=0.3em]
   \item Students are put into small groups and are asked to write down their confidence on a single notecard without names.
   \item The cards are evidence of the group's confidence in the content.
\end{enumerate}
\\ \hline
\end{tabular}

\vspace{0.5cm}
\noindent
{\large\textbf{Content (fully worked out problems \& answers)}

\begin{tabular}{|p{0.9\linewidth}|}
\hline
\begin{itemize}
   \item No new worked problems; brain dump.
\end{itemize}
\\ \hline
\end{tabular}

\vspace{0.5cm}
\noindent
{\large\textbf{Checking for Understanding (questions \& answers) continued...}

\begin{tabular}{|p{0.9\linewidth}|}
\hline
\begin{itemize}
   \item Review confidence levels
   \item Interpret the variety of confidence and their distribution in groups.
   \item Get feedback for what they like and don't like
   \item Check in if the time is suitable for the students
\end{itemize}
\\ \hline
\end{tabular}
\newpage
\section*{Plan Checklist}

\begin{tabular}{|p{0.9\linewidth}|}
\hline
\begin{itemize}
\item[$\Box$] Variety of Collaborative Learning Techniques and Learning Strategies
\item[$\Box$] Chapter/slide\# where information is coming from
\item[$\Box$] Included timestamps (and buffer time as needed)
\item[$\Box$] Details on how leader will break up students
\item[$\Box$] Checking for Understanding questions
\item[$\Box$] Fully worked out problems and examples for each activity
\item[$\Box$] Necessary links required for online implementation (if applicable)
\item[$\Box$] Source material correctly cited (if applicable)
\item[$\Box$] Plan is uploaded to Teams
\end{itemize}
\\ \hline
\end{tabular}
\newpage

\end{document}